\chapter{Summary, Conclusions, and Future Work}
\label{chap:conclusion}
\chapterrunning{summary and conclusions}

The dissertation addressed a narrow question: can an LLM-guided code-evolution loop generate
useful heuristic adaptation under controlled evaluation? The completed evidence supports a
limited positive answer. Across several families, the loop produces executable diversity and sometimes
preserves useful refinements. The same evidence also shows that the resulting gains are strongly
bounded by validator pressure, baseline strength, benchmark headroom, and the exact search object
under evolution.

\section{Cross-Family Synthesis}
\label{sec:conclusion-synthesis}

The completed phases reveal recurring patterns rather than one isolated success or failure. Those
patterns are best summarized at two scales. Table~\ref{tab:conclusion-atlas} is a master evidence
atlas of the retained empirical record. Table~\ref{tab:conclusion-patterns} then distills the
recurring cross-family patterns that survive that evidence.

\clearpage
\begin{landscape}
\begingroup
\scriptsize
\setlength{\LTleft}{0pt}
\setlength{\LTright}{0pt}
\setlength{\tabcolsep}{1.5pt}
\renewcommand{\arraystretch}{1.04}
\begin{longtable}{
  >{\raggedright\arraybackslash}p{0.33in}
  >{\raggedright\arraybackslash}p{0.55in}
  >{\raggedright\arraybackslash}p{0.65in}
  >{\raggedright\arraybackslash}p{0.67in}
  >{\raggedright\arraybackslash}p{0.52in}
  >{\raggedright\arraybackslash}p{0.66in}
  >{\raggedright\arraybackslash}p{0.86in}
  >{\raggedright\arraybackslash}p{0.75in}
  >{\raggedright\arraybackslash}p{0.55in}
  >{\raggedright\arraybackslash}p{0.70in}
  >{\raggedright\arraybackslash}p{1.16in}
}
\caption[Master evidence atlas]{Master evidence atlas across the completed dissertation phases. ``Baseline / control'' names the
phase-specific comparator used for the reported main delta. Lower deltas are better. Cells are
marked ``---'' when the phase does not admit one matched paired contrast without forcing an
artificial comparison.}
\label{tab:conclusion-atlas}\\
\toprule
Phase & Family & Search object & Conditions & N / pairing & Primary endpoint & Best LLM condition & Baseline / control & Main delta & 95\% CI & Interpretation \\
\midrule
\endfirsthead
\multicolumn{11}{l}{\small\itshape Table \thetable\ continued from the previous page}\\
\toprule
Phase & Family & Search object & Conditions & N / pairing & Primary endpoint & Best LLM condition & Baseline / control & Main delta & 95\% CI & Interpretation \\
\midrule
\endhead
\midrule
\multicolumn{11}{r}{\small\itshape Continued on the next page}\\
\endfoot
\bottomrule
\endlastfoot
1--4 & simple games & full policy code & 3 envs,\newline 3 reviewed recipes & 10 transfer runs & held-out win rate,\newline margin & territory control,\newline win rate $0.9467$,\newline margin $34.17$ & --- & --- & --- & Held-out transfer is measurable, but novelty spikes still collapse into churn often enough that code novelty alone is not evidence of adaptation. \\
5 & TSP & bounded heuristic code & 4 replay arms & 20 paired offsets & held-out TSPLIB gap & random replay,\newline mean $0.152211$ & no replay,\newline mean $0.183325$ & $-0.0311$ & $[-0.0602,$\newline$-0.0008]$ & Clearest replay win in the thesis; the best arm is random replay rather than raw failure replay. \\
5 & ATSP & bounded heuristic code & 4 replay arms & 20 paired offsets & held-out ATSP gap & compressed failure replay,\newline mean $0.187225$ & no replay,\newline mean $0.192847$ & $-0.0056$ & $[-0.0191,$\newline$0.0071]$ & Compression-aware replay is mildly favorable, but the primary endpoint is not cleanly separated under paired uncertainty. \\
5 & CVRP & bounded heuristic code & 4 replay arms & 20 paired offsets & held-out CVRPLIB gap & no replay,\newline mean $0.235779$ & random replay,\newline mean $0.246032$ & $-0.0076$ & $[-0.0146,$\newline$-0.0006]$ & Replay is not a stable default even within routing; the strongest family-level condition in the replicated extension is the no-replay arm. \\
6 & TSP & bounded heuristic code & 7 replay arms & 20 paired offsets & held-out TSPLIB gap & diversity-aware failure replay & raw failure replay & $-0.008637$ & $[-0.035972,$\newline$0.017862]$ & Replay curation matters more than archive size alone; archive hardness tracks worse held-out transfer more strongly than descriptor diversity does. \\
7 & TSP & modular operator code & 7 conditions & 20 paired offsets & held-out TSPLIB gap,\newline transplant,\newline ablation & full-solver evolution,\newline TSPLIB $0.164668$ & heuristic-only baseline & $-0.003035$ & $[-0.010147,$\newline$0.000716]$ & Rediscovery signatures recur, but no operator survives the full validation pipeline strongly enough to support a reusable-operator claim. \\
8 & TSP & controller over frozen portfolio & 8 conditions & 20 paired offsets & held-out TSPLIB gap,\newline selector regret & adaptive controller,\newline mean $0.155397$ & static LLM selector & $-0.038118$ & $[-0.06868,$\newline$-0.00707]$ & Adaptive control beats a weak one-shot selector, but the fixed heuristic remains best overall because oracle headroom is near zero on the primary endpoint. \\
9 & CVRP & bounded whole-solver code & 4 conditions & 20 paired offsets & held-out penalized gap & solver evolution,\newline mean $0.182231$ & Clarke--Wright,\newline mean $0.073766$ & $+0.108465$ & $[0.071184,$\newline$0.144731]$ & Feasible solver evolution beats weaker baselines, but it does not surpass the strongest fixed heuristic and it carries a heavy runtime tail. \\
9 closeout & CVRP & bounded whole-solver code & 6 conditions,\newline matched learned controls & 20 paired offsets & held-out penalized gap & replay solver evolution,\newline mean $0.171114$ & budget-matched no replay,\newline mean $2.932746$ & $-2.761632$ & $[-3.369777,$\newline$-2.201586]$ & A replay-specific advantage remains after matched learned controls are added, but the strongest fixed Clarke--Wright baseline still remains best overall. \\
\end{longtable}
\endgroup
\end{landscape}
\clearpage

The atlas separates three things that are easy to conflate in a code-evolution project: whether the
loop can generate executable diversity at all, whether that diversity survives held-out evaluation,
and whether the resulting gain remains visible once the strongest available baseline is used as the
reference point.

In the archival phase chronology, the atlas covers the retained evidence from phases 1 through 10
in bundled form, together with the scoped phase-9 closeout on the same bounded CVRP regime. The
first row aggregates the simple-game record from phases 1, 2, 2b, 3, and 4; phases 5 through 9
appear directly; the closeout records the final matched-budget mechanism study on real-world CVRP;
and phase 10 is the synthesis step used to interpret those retained results rather than a new
benchmark family. The scoped phase-11 model-strength factorial remains outside the retained
dissertation evidence base.

\begin{table}[t]
\centering
\scriptsize
\setlength{\tabcolsep}{2pt}
\begin{tabularx}{\textwidth}{
  >{\raggedright\arraybackslash}p{0.95in}
  Y
  Y
  Y
}
\toprule
Evidence family & Strongest positive signal & Main limiting factor & Dominant failure mode \\
\midrule
Simple games & measured behavioral transfer, especially in territory control & environment dependence & lexical novelty without robust behavioral gain \\
Routing replay (phases 5--6) & replay helps on symmetric TSP & family-specific mechanism sensitivity & unstable transfer of replay recipes across TSP, ATSP, and CVRP \\
Operator discovery (phase 7) & recurring operator rediscovery signatures & strict transplant and ablation criteria & scaffold-specific or brittle operators \\
Adaptive portfolio (phase 8) & adaptive controller beats weak one-shot LLM selector & near-zero oracle headroom on the primary endpoint & lack of complementarity in the frozen portfolio \\
Real-world CVRP (phase 9) & feasible solver evolution beats weaker baselines & strong Clarke--Wright baseline and runtime variability & bounded improvement without strongest-baseline dominance \\
Real-world CVRP closeout (phase 9 closeout) & replay-aware search beats both matched learned controls & strong Clarke--Wright baseline still remains best overall & memoryless or direct learned controls collapse on feasibility and gap \\
\bottomrule
\end{tabularx}
\caption[Recurring cross-family patterns]{Recurring patterns across the completed phases. The positive signals are real, but each is
paired with a concrete limiting factor rather than with an unrestricted win narrative.}
\label{tab:conclusion-patterns}
\end{table}

Four synthesis statements follow. Syntactic code novelty is easy for the loop to produce, but
novelty is not sufficient for functional gain. Improvements are easiest to obtain when the
benchmark leaves reachable headroom and when the baseline is not already very strong. The dominant
failure mode changes by family: churn in games, replay-selection sensitivity in routing,
validation collapse for modular discovery, selector-headroom limits in portfolios, and
feasibility/runtime pressure in real-world CVRP. The most conservative interpretation consistent
with the retained evidence is capability-bounded adaptation rather than general algorithm
discovery.

The direct single-shot Codex CVRP baseline reinforces this final point. Evaluated descriptively
through the same phase-9 validator, the interactive Codex-authored solver reaches full held-out
feasibility with mean penalized gap $0.06778$, which is close to the Clarke--Wright baseline and
well ahead of both the replicated autonomous solver-evolution mean and the bounded direct closeout
control. That comparison is descriptive only, not a matched experimental condition, but together
with the closeout it is consistent with the possibility that current interactive coding assistance
can extract more stable value from the model in this regime than the bounded autonomous protocols
retained in the dissertation.

\section{Assessment of the Working Hypotheses}
\label{sec:conclusion-hypotheses}

Table~\ref{tab:conclusion-hypotheses} revisits the hypotheses stated in
Section~\ref{sec:intro-hypotheses}.

\begin{table}[t]
\centering
\scriptsize
\setlength{\tabcolsep}{3pt}
\begin{tabularx}{\textwidth}{
  >{\raggedright\arraybackslash}p{0.52in}
  >{\raggedright\arraybackslash}p{0.9in}
  Y
}
\toprule
Hyp. & Status & Evidence-based assessment \\
\midrule
H1 & Supported & The loop generates executable diversity and produces useful held-out adaptation in the simple games, positive replay transfer on symmetric TSP, and feasible bounded improvement on real-world CVRP. \\
H2 & Partially supported & Replay helps in some settings, and the phase-9 closeout shows a clear replay-specific advantage over matched learned controls. Even so, no single replay recipe dominates across the retained families, and the strongest fixed baselines still prevent a universal replay win claim. \\
H3 & Supported & As validators strengthen and classical baselines consume more benchmark structure, the gains narrow. This pattern is strongest in phases 8 and 9, where the best fixed or supervised baselines remain ahead on the primary endpoint. \\
H4 & Supported & Transfer exists, but it is partial and family-dependent. The strongest early transfer is concentrated in territory control, and the routing phases show that signals learned on one family do not carry over uniformly to others. \\
\bottomrule
\end{tabularx}
\caption[Assessment of the working hypotheses]{Assessment of the working hypotheses. The evidence supports adaptation, but in a bounded
and family-sensitive form.}
\label{tab:conclusion-hypotheses}
\end{table}

This hypothesis-level summary captures the main contribution of the dissertation. The project does
not end with a single dominant heuristic recipe. Instead, it leaves a more precise empirical
account of where the loop transfers, where it weakens, and what kinds of controls are needed to
separate those cases.

\section{Limitations}
\label{sec:conclusion-limitations}

Several limitations should remain explicit.

\begin{enumerate}[leftmargin=2em]
\item The routing phases rely on bounded scaffolds and interfaces. This choice improves
interpretability, but it also narrows the search space relative to unconstrained program synthesis.
\item The completed campaigns use fixed benchmark subsets rather than the full diversity of every
public benchmark library.
\item The direct Codex CVRP baseline is descriptive and human-guided rather than replicated and
randomized.
\item The dissertation does not retain a completed crossed experiment that cleanly separates base
model tier from evolution technique under matched budgets. The phase-9 closeout resolves the
replay-versus-budget question within one fixed model tier, but cross-tier attribution remains open.
\item Some family-level comparisons are qualitative by necessity because the primary metrics differ
across games, TSP/ATSP, and CVRP.
\end{enumerate}

These limitations do not invalidate the central claims. They define the boundary of those claims
and help explain why the final interpretation is intentionally conservative.

\section{Future Work}
\label{sec:conclusion-future-work}

One possible next step is a clean model-strength continuum study. A suitable design would cross task
family, model tier, and evolution technique under matched candidate budgets, with at least
single-shot, no-replay, and replay-based arms. Such a study would answer whether the gains
observed here come mainly from the base model, from the outer search loop, or from their
interaction.

Three additional directions follow directly from the retained results.

\begin{enumerate}[leftmargin=2em]
\item Replay should be revisited through curation rather than through larger archives alone. Phase 6
indicates that hardness and coverage matter more than simple diversity counts.
\item Adaptive portfolio work should begin with stronger portfolio complementarity. If the oracle
selector cannot beat the single best heuristic on the primary endpoint, then controller learning has
little room to show a clear gain.
\item Real-world solver evolution should be extended to harder routing families, stronger external
references, and explicit runtime-quality trade-off analysis. CVRPTW and broader CVRPLIB regimes are
natural candidates once the bounded X-instance setup is exhausted.
\end{enumerate}

A further methodological direction is to formalize the direct interactive-coding baseline. The
current Codex-authored CVRP solver is useful descriptively, but a frozen and replicated human-in-the-loop
protocol would allow a fairer comparison between autonomous and interactive use of the same model
class.

\section{Final Conclusion}
\label{sec:conclusion-final}

LLM-guided code evolution is scientifically worth studying because it can generate executable
heuristic variation and, in some settings, preserve useful adaptation across nontrivial domains.
At the same time, the completed evidence does not support a claim of broad autonomous solver
discovery. The dissertation's main contribution is a cross-family methodology and evidence base
for distinguishing useful heuristic evolution from code churn, and bounded progress from
overstatement. That is a narrower claim than early enthusiasm around LLM-driven algorithm
discovery often suggests, but it is the stronger and more defensible scientific result.
